% !PRINTED 202
\weberpage{202}

From this it now follows that the division by 2, and hence also by
every power of 2, can be carried out by a chain of square roots. If we
therefore assume the problem of the division by an odd number to be
solved, then the division by an even number is reduced to
square roots. In what follows we shall occupy ourselves exclusively
with \gesperrt{the division by odd numbers}.

\begin{center}
§ 60. \gesperrt{The Division by an Odd Number.}
\end{center}

If we put, when $n$ is an odd number,

\weq{1}{x = \mathrm{sn}\,\frac{v}{n}, \qquad y = \mathrm{cn}\,\frac{v}{n}, \qquad z = \mathrm{dn}\,\frac{v}{n},}

then from the equations IV, \S\ 57 we obtain:

\weq{2}{\begin{aligned}
D(x^2)\,\mathrm{sn}\,v - x\,A(x^2) &= 0,\\
D(x^2)\,\mathrm{cn}\,v - y\,B(x^2) &= 0,\\
D(x^2)\,\mathrm{dn}\,v - z\,C(x^2) &= 0.
\end{aligned}}

The first of these equations is of degree $n^2$ with respect to the
unknown $x$; by the second and the third, $y$ and $z$ are expressed
\gesperrt{rationally} in terms of $x$ (and of cn\,$v$, dn\,$v$). There result,
then, $n^2$ systems of values for the three unknowns $x, y, z$, whose
meaning is:

\weq{3}{\begin{aligned}
x_{\mu,\mu'} &= \mathrm{sn}\left(\frac{v}{n} + \frac{4\mu K+4\mu' iK'}{n}\right),\\
y_{\mu,\mu'} &= \mathrm{cn}\left(\frac{v}{n} + \frac{4\mu K+4\mu' iK'}{n}\right),\\
z_{\mu,\mu'} &= \mathrm{dn}\left(\frac{v}{n} + \frac{4\mu K+4\mu' iK'}{n}\right),
\end{aligned}}

in which $\mu, \mu'$ each run through a complete system of residues
$(\mathrm{mod}\ n)$.

The first equation (2):

\weq{4}{D(x^2)\,\mathrm{sn}\,v - x\,A(x^2) = 0,}

of degree $n^2$, is called the \gesperrt{general division equation}. It
is irreducible in the sense that it is not decomposable into factors
which are rational with respect to sn\,$v$, cn\,$v$, dn\,$v$ and have
arbitrary coefficients independent of $v$.

For, in the first place, the $n^2$ values $x_{\mu,\mu'}$ are all
different from one another, and if any rational equation:

\[
F\left(\mathrm{sn}\,\frac{v}{n},\ \mathrm{sn}\,v,\ \mathrm{cn}\,v,\ \mathrm{dn}\,v\right) = 0
\]
%CONT>>
